\documentclass[a4paper]{article}
\usepackage{graphicx}
\usepackage{amsmath,amssymb}
\usepackage{hyperref}
\usepackage{minted}

\title{Crazy fly}
\date{}
\begin{document}
    \maketitle
    
    
    \section{Crazyflie 2.1+ Overview}
        
        \textbf{Platform:} Crazyflie 2.1+ is a palm-sized open-source micro quadcopter designed primarily for research and development rather than consumer applications.
        
        \subsection{Physical Characteristics}
            
            \begin{itemize}
                \item Weight: approximately 29 grams
                \item Motor-to-motor distance: approximately 92 mm
                \item Fully palm-sized form factor
            \end{itemize}
        
        \subsection{Onboard Sensors (IMU)}
            
            The main board includes:
            
            \begin{itemize}
                \item 3-axis accelerometer and 3-axis gyroscope (BMI088)
                \item 3-axis magnetometer
                \item Barometer (BMP388)
            \end{itemize}
            
            These sensors enable state estimation, attitude control, and onboard filtering (e.g., EKF).
        
        \subsection{Processing Units}
            
            \begin{itemize}
                \item STM32F405 (Cortex-M4) for flight control
                \item nRF51822 for radio communication
            \end{itemize}
        
        \subsection{Communication Interfaces}
            
            \begin{itemize}
                \item 2.4 GHz radio communication (Crazyradio)
                \item USB interface
                \item Optional WiFi via AI Deck
            \end{itemize}
        
        \subsection{Expansion Capability (Deck System)}
            
            Crazyflie uses a modular expansion deck system, allowing additional hardware to be mounted directly on top of the platform.
            
            \subsubsection{Multi-Ranger Deck}
            
            \begin{itemize}
                \item Five Time-of-Flight (ToF) laser sensors
                \item Measures distance in front, back, left, right, and upward directions
                \item Range up to approximately 4 meters
                \item Functions as a lightweight multi-directional ranging system
                \end{itemize}
            
            \subsubsection{Flow Deck v2}
            
            \begin{itemize}
                \item Optical flow sensor
                \item Downward-facing ToF sensor for altitude measurement
                \item Enables stable indoor flight without GPS
                \end{itemize}
            
            \subsubsection{AI Deck 1.1}
            
            \begin{itemize}
                \item GAP8 RISC-V processor
                \item 320×320 monochrome camera (Himax HM01B0)
                \item ESP32 for WiFi connectivity
                \item Supports onboard AI workloads and custom vision processing
                \end{itemize}
        
        \subsection{Software Characteristics}
            
            \begin{itemize}
                \item Fully open-source firmware
                \item Python API available
                \item Firmware-level modification possible
                \item Compatible with ROS
                \item Simulation support (e.g., Gazebo)
            \end{itemize}
        
        \subsection{Advantages}
            
            \begin{itemize}
                \item Extremely lightweight and compact
                \item Fully programmable at firmware level
                \item Modular sensor expansion
                \item Suitable for swarm research
                \item Suitable for indoor autonomy experiments
            \end{itemize}
        
        \subsection{Limitations}
            
            \begin{itemize}
                \item Very limited payload capacity
                \item Short flight time (approximately 7–10 minutes)
                \item Poor wind resistance
                \item Ranging sensors are ToF-based rather than scanning LiDAR
            \end{itemize}
        
        \subsection{Conclusion}
            
            For a minimal fully-programmable quadcopter platform with onboard IMU, expandable ranging sensors, and optional camera-based AI processing, Crazyflie 2.1+ with Multi-Ranger and AI Deck represents one of the smallest practical research platforms currently available.

\end{document}